% Generated by analyze_sparse_architecture_ablation.py.
\begin{table}[!htbp]
\centering
\caption{Coupled routed operating points for the canonical 17-tile request. Time is cycles divided by each row's realized clock; it is not a matched-build speedup. Resource and memory differences prevent isolated marginal attribution.}
\label{tab:organization-operating-points}
\small
\setlength{\tabcolsep}{2.6pt}
\begin{tabular}{lrrrrrrrr}
\toprule
Org. & Mcyc. & MHz & ms & AXI MB & LUT & FF & BRAM & DSP \\ \midrule
Independent & 9.858 & 100.0 & 98.577 & 2.400 & 38,274 & 28,681 & 52.0 & 193 \\
Rectangle & 7.432 & 100.0 & 74.321 & 2.386 & 38,274 & 28,681 & 52.0 & 193 \\
Persistent & 4.923 & 50.0 & 98.458 & 1.322 & 44,408 & 26,056 & 3.5 & 180 \\
+ commit & 5.134 & 50.0 & 102.676 & 1.322 & 45,046 & 26,648 & 3.5 & 180 \\
\bottomrule
\end{tabular}
\end{table}
\noindent\textbf{Separate 50-MHz builds:} Mask-only and commit-safe DMA-trigger-to-completion means are 102.677 and 102.677 ms. The latter build uses an additional 638 LUTs and 592 registers, with no additional DSP or BRAM. This is not an isolated publication-overhead measurement or complete application latency. Cycle counters use different terminal boundaries.
